\chapter{Die allumfassende Architektur}
\label{chap:die-allumfassende-architektur}

\section{Auswertung vorhandener Produkte}
% Qlik, Tableau, Grafana => Vergleichem

\section{Workflow of a Data Scientist}
% Collection Exploration Mugling Modeling Validation Reporting

\section{Mögliche Veranschaulichungen von Daten}
% Aktuelle Modelle: Measures and Dimensions

\section{Mircoservices}

\subsection{Microfrontend}
Domain Driven Development also in Frontend.

Aufteilung in Module, die während der laufzeit geladen werden.
Seite 171 in Buch (Eberhard Wolff Microservices) befasst sich mit dem
Thema, wenn auch nicht auf das dynamische Laden eingegangen wird.
"Das Deployment der SPA ist meistens nur als vollständige Anwendung möglich."
Neuer Ansatz für unabhängiges Deployment!

-> Auch Problem: Deployment Monolith -> Separation der Deployment Vorgänge

-> "Wenn in einem Microservice ein Feature umgesetzt wird, das auch Änderungen
in der Client-Anwendung benötigt, kann diese Änderung nicht durch eine neue
Version des Microservice alleine ausgerollt werden. Es muss auch eine neue Version
der Client-Anwendung ausgeliefert werden." Seite 176.  (Eberhard Wolff Microservices)
Absatz 3.

Deployment abhängigkeiten verringern!

SOAP Schnittstelle zwischen Frontend Chart Plugins und Frontend Arrangement Software!
Weniger bis garkeine Abhängigkeit von Implementierungsdetails

Problem: Steigende Komplexität der Anwendung.

\section{Die Schnittstellengestaltung}
% sehr wichtig, abstraktion

\section{Die Aufteilung der Logik}
\label{sec:die-aufteilung-der-logic}
Separation of Concerns, etwa Trennung der Zuständigkeiten

\section{Die Datenstruktur}
Die Softwarearchitektur sollte so flexibel wie möglich mit unterschiedlichen Datenstrukturen
umgehen können. Dennoch benötigt es eine klar definierte Schnittstelle zu den Plugins,
welche für die Datenvisualisierung verantwortlich sind. Aber wie sollte nun eine solche flexible 
und dennoch klar definierte Datenstruktur aussehen?

Schauen wir uns zuerst einmal die bereits existierenden Lösungen verwandter Anwendungen an. Drei der
wohl bekanntesten Business Intelligence Lösungen im Bereich der Datenvisualisierung sind Qlik Sense
von Qliktech, Tableau Public von Tablea Software und PowerBI von Microsoft. In der Business Intelligence
Welt werden Mengen von Elementen oft in Maße (Measures) und Dimensionen (Dimensions) unterteilt. Unter Maße versteht
man numerische Elemente einer Menge, Dimensionen hingegen stellen kategorische Elemente dar. Ein Maß
ist beispielsweise die Anzahl verkaufter Produkte, eine Dimension das Land, in welchem die Produkte verkauft worden
sind. Bei einem Balkendiagram kann demnach die Länge der Balken die Anzahl der verkauften Produkte
widerspiegeln. Ein Balken representiert das Land, in dem das Produkt verkauft wurde.



Die Abbildung von Mehrdimenstionaler Daten in einer eindimensionalen Welt.

% zwei möglich eoperationen nötig, alle beziehungen in einer tabelle darstellen,
% zeilen unendliche Teilmengen hinzufügen und die Sortieren
% List processing???

Da die Daten und deren Beziehungen von Anwendungsfall zu Anwendungsfall
variieren, gilt es eine klar definierte Datenstruktur zu finden, die
die gängigsten Anwendungsfälle abdecken kann. 

% Show dashboard json structure using maybe css grid
% and data structure to pass to plugins
% many to many to many relations???



In Tableau und Qlik Sense hat man die Möglichkeit, aus einer ausklappbaren Auswahl an der Seite der
Anwendung, Maße und Dimensionen durch Drag and Drop verschiedenen Diagrammen zuzuordnen. 

\section{Das Design des Datenflusses}
\label{sec:das-design-des-datenflusses}

\section{Die Definition der Schnittstellen}
\label{sec:die-definition-der-schnittstellen}

\section{Die Wahl der Technologien}

\subsection{React}
Dan Abramov Representational components vs container components => actual view and controller


\section{Die Bewahrung der Offenheit gegenüber der Veränderung}
\label{sec:die-bewahrung-der-offenheit-gegenueber-der-veraenderung}

\subsection{Die Änderung des Quellcodes}

\subsection{Das stetige Aktualisieren der ausgelieferten Software}
Probleme: cookies verändern sich, serviceworker verändert sich, neue Felder in der Datenbank
Frontend und Backend etc, also allen verschiedenen Services


\begin{listing}
    \label{lst:HelloJSX}
    \caption{Ein einfaches JSX Beispiel}
    \inputminted{jsx}{snippets/examples/Welcome.jsx}
\end{listing}

\begin{listing}
    \label{lst:Golang}
    \caption{Ein einfaches Golang Beispiel}
    \inputminted{go}{snippets/examples/hello.go}
\end{listing}

\newpage

\begin{figure}
    \label{figure:test}
    \includegraphics[scale=0.2]{img/HTW}
    \caption{Test}
\end{figure}

\begin{figure}
    \label{figure:beispiel}
    \includegraphics[scale=0.2]{img/HTW}
    \caption{Beispiel}
\end{figure}
